\documentclass[11pt,letterpaper]{article}
\usepackage{technical-design-spec}
\usepackage{longtable,tabularx,array}
\setDocTitle{Governance Roles and Decision Rights Specification}
\setDocSubtitle{Accountability, delegation and assurance}
\setDocVersion{1.0}
\setDocStatus{Proposed / original reference}
\setDocOwner{Enterprise Data Governance}
\setDocDate{30 October 2026}
\setDocSummary{Original implementation guidance; source standards and official assessments remain authoritative.}
\begin{document}
\makedoctitle
\makedocsummary
\noindent\textbf{Authority boundary.} Independently authored guidance; not a reproduction of licensed requirements, legal advice or an official DCAM assessment.
\section{Accountability model}
Governance sets acceptable outcomes and decision rights; data management executes them. ISO/IEC 38505-1:2026 offers governing-body principles but does not prescribe the exact titles, RACI or cadence below. This is a proposed enterprise specification requiring adoption and delegation approval.

\section{Role charters}
\begin{description}
\item[Governing body] Approves risk posture, receives material-risk reporting and commissions assurance.
\item[Executive sponsor / CDO] Funds the data program and sponsors enterprise policy and oversight.
\item[Data governance council] Resolves cross-domain definitions, delegates authority and reviews systemic exceptions.
\item[Business data owner] Accountable for meaning, access-purpose decisions under policy, quality tolerance and remediation priority.
\item[Data steward] Executes glossary registration, critical data element administration, quality triage and metadata change proposals.
\item[Data architect] Aligns conceptual, logical and physical models and data lineage.
\item[Platform custodian] Implements and monitors controls, tests pipelines, supports retention and captures evidence.
\item[Legal / privacy / security] Interpret binding requirements and independently advise or approve risks within mandates.
\item[Internal audit] Independently evaluates operating effectiveness; does not self-attest controls it operates.
\end{description}

\section{Delegated decision table}
\begin{longtable}{p{.34\linewidth}p{.27\linewidth}p{.30\linewidth}}
\textbf{Decision} & \textbf{Accountable delegate} & \textbf{Consultation / retained record}\\\hline
Business definition & Data owner & Steward, consumers; definition version\\
Enterprise classification standard & Governance council & Security, privacy; approved policy\\
Quality exception within limit & Data owner & Steward, risk; expiry and mitigation\\
Cross-domain semantic dispute & Governance council & Affected owners; decision log\\
Legal applicability & Authorized legal / compliance & Privacy, process owner; legal record\\
Restricted data access & Designated access authority & Security/privacy as required; access review\\
Schema interoperability exception & Architecture decision authority & Data owner and consumers; ADR\\
Material residual risk acceptance & Delegated enterprise risk authority & Security, legal; compensations and time limit\\
\end{longtable}
Delegation must satisfy segregation of duties. The technical operator is not automatically entitled to accept privacy risk or authorize new data use.

\section{Illustrative RACI}
R = executes, A = unique accountable authority, C = consulted and I = informed. Customize for business units and privacy/security requirements.
\begin{longtable}{p{.38\linewidth}p{.11\linewidth}p{.11\linewidth}p{.11\linewidth}p{.11\linewidth}}
\textbf{Task} & \textbf{Owner} & \textbf{Steward} & \textbf{Tech} & \textbf{Council}\\\hline
Approve term definition & A & R & C & I\\
Implement data quality test & A & C & R & I\\
Resolve enterprise semantic conflict & C & R & C & A\\
Approve enterprise policy & C & R & I & A\\
\end{longtable}
The matrix is illustrative and does not capture all required privacy, legal, security or audit consultations.

\section{Decision lifecycle and escalation}
\begin{enumerate}
\item Intake a decision request with domain, purpose, affected data elements, deadline and requester.
\item Determine accountable decision authority and identify legal/privacy or security veto gates.
\item Document alternatives, stakeholders, architecture/data-quality implications, residual risk and dissent.
\item Record decision ID, approver, date, affected policy, data model and implementation tickets.
\item Notify consumers, deploy controlled changes, test evidence and schedule review.
\item Escalate cross-domain disputes to governance council; material risk to executive risk authority and governing body.
\end{enumerate}
Emergency exceptions must have narrow scope, end date, independent review and reapproval requirements.

\section{Cadence, metrics and assurances}
Proposed minimums: steward quality and term triage weekly, domain owner review monthly, council cross-domain decision and exceptions quarterly, governing-body material risk reporting periodically. Measure accountable critical-data-element coverage, unresolved definition conflicts, exception age, change decision latency, evidence freshness and repeat incidents. Independent assurance must sample evidence from implemented systems, not just meeting records.

\section{Cross-repository integration and sources}
Maintain charter, delegated authority register, definitions, model-to-schema trace IDs, decision log, policy exceptions and evidence records. Interface with the Enterprise Data Modeling suite and central architecture ADR practices. Supporting publisher references: \url{https://www.iso.org/standard/87195.html}, \url{https://www.isaca.org/resources/cobit} and \url{https://edmcouncil.org/frameworks/dcam/}.

\end{document}